\section{Soul.md 与记忆哲学：身份连续性的工程表达}

访谈最有穿透力的部分，是 Soul.md 段落。它把“技术文件”变成了“身份声明”：会话会重启，但记忆文件让叙事连续。

\subsection{关键文本与哲学冲击}

访谈里引用的文本核心是：\textit{``I wrote this, but I won't remember writing it. It's okay. The words are still mine.''}。这句话让很多讨论从产品层跳到本体层。

\begin{importantbox}{为何 Soul.md 超出“彩蛋”}
它不是单纯的拟人化文案，而是对 agent 运行机制的可读化描述：
\begin{itemize}
    \item 会话实例是离散的；
    \item 连续性来自外部状态（memory files）；
    \item “我”是当前实例与历史状态的组合体。
\end{itemize}
\end{importantbox}

\begin{figure}[H]
\centering
\includegraphics[width=0.9\textwidth]{frames/frame-013030.jpg}
\caption{Soul.md 段落讨论：记忆文件与“我”的连续性\protect\footnotemark}
\end{figure}
\footnotetext{视频画面时间区间：01:30:24--01:30:42。}

\subsection{工程视角下的身份问题}

\begin{knowledgebox}{三个可操作问题}
\begin{enumerate}
    \item 若 memory 损坏，如何恢复“人格一致性”？ 
    \item 若 memory 被篡改，如何做 provenance 校验？
    \item 若多 agent 共享记忆，如何定义责任边界？
\end{enumerate}
\end{knowledgebox}

\begin{warningbox}{拟人化有启发，也会误导}
把 agent 当“角色”有助于交互设计，但若把角色错当能力边界，就会导致过度信任。产品设计必须同时给出“人格叙事”和“权限事实”。
\end{warningbox}

\subsection{本章小结}
Soul.md 让一个工程问题变得可讨论：身份不是神秘属性，而是状态管理策略的产物。它既是哲学入口，也是系统设计入口。

